‫% -------------------------------------------------------
‫%  Abstract
‫% -------------------------------------------------------
‫
‫\section{مقدمه}
‫
‫هدف این بخش، مشاهده مستقیم سازوکار بارگذاری یک صفحه وب مبتنی بر \لر{HTTP} و پاسخ‌گویی
‫به پرسش‌های آماری و تحلیلی صورت آزمایش است. ضبط روی \لر{Wireless Interface}
‫\لر{wlp0s20f3} انجام شد.
‫همان‌گونه که در شکل~\لر{\رجوع{شکل:واسط-ضبط}} دیده می‌شود، این \لر{Interface} در صفحه آغازین وایرشارک
‫فعال و دارای ترافیک بود.
‫
‫در بخش دوم، هدف مشاهده ارتباط متنی دوطرفه \لر{Telnet}، \لر{Telnet Option Negotiation}
‫و میزان خوانایی داده‌های \لر{Client} و سرور در فایل ضبط است. برای این منظور،
‫نشست\پانویس{Session} زنده \لر{telehack.com} و شواهد
‫ارائه‌شده از \لر{telnet.pcap} به‌صورت جداگانه بررسی شدند.
‫
‫در بخش سوم، پس از پاک‌کردن حافظه نهان \لر{DNS}، یک پرس‌وجوی نوع \لر{A} با
‫\لر{nslookup} ارسال و زوج درخواست و پاسخ آن در وایرشارک تحلیل شد. هدف، تعیین سرور مقصد و
‫تفسیر دقیق سرآیند، شمارنده‌ها و رکورد پاسخ \لر{DNS} بود.
‫
‫\شروع{شکل}[H]
‫\centerimg{1.png}{14cm}
‫\شرح{انتخاب \لر{Interface} فعال \لر{wlp0s20f3} در \لر{Wireshark 4.4.8}}
‫\برچسب{شکل:واسط-ضبط}
‫\پایان{شکل}
‫
‫پس از آغاز ضبط، ورود بسته‌ها در فهرست اصلی مشاهده شد. شکل~\لر{\رجوع{شکل:ضبط-زنده}} نمایی از
‫ضبط زنده را نشان می‌دهد. فایل نهایی \لر{http.pcap} شامل \لر{5953} بسته، \لر{2074855} بایت داده و
‫\لر{82.379578751} ثانیه ترافیک است.
‫
‫\شروع{شکل}[H]
‫\centerimg{2.png}{14cm}
‫\شرح{ضبط زنده ترافیک روی \لر{Wireless Interface}}
‫\برچسب{شکل:ضبط-زنده}
‫\پایان{شکل}
‫
‫\section{تقسیم کار}
‫
‫آزمایش اوّل توسط کیارش شجاعی و آزمایش دوّم و سوّم توسط آریانا زال‌نژاد انجام شد. هرکس گزارش‌ و تصاویر خود را تهیه کرد و گزارش نهایی توسط کیارش شجاعی جمع‌آوری شد.
‫
‫\section{روش انجام آزمایش}
‫
‫نشانی \href{http://sina.sharif.edu/\%7Eazad/}{\لر{http://sina.sharif.edu/\textasciitilde{}azad/}}
‫در پنجره ناشناس مرورگر \لر{Chrome} باز شد.
‫نمای کامل صفحه و استفاده از \لر{HTTP} ناامن در شکل~\لر{\رجوع{شکل:صفحه-وب}} دیده می‌شود.
‫پس از بارگذاری و بازخوانی صفحه، ضبط متوقف و فایل با قالب \لر{pcap} ذخیره شد. نشانی رایانه
‫آزمایش \لر{10.22.10.57} و نشانی سرور وب \لر{152.89.13.54} بود.
‫
‫\شروع{شکل}[H]
‫\centerimg{3.png}{14cm}
‫\شرح{صفحه وب هدف در پنجره ناشناس مرورگر}
‫\برچسب{شکل:صفحه-وب}
‫\پایان{شکل}
‫
‫برای یافتن درخواست‌ها، فیلتر \لر{http.request.method == "GET"} اعمال شد. مطابق
‫شکل~\لر{\رجوع{شکل:درخواست‌های-get}}، نخستین درخواست صفحه در فریم \لر{3541} از
‫\لر{10.22.10.57} به \لر{152.89.13.54} ارسال شده و سرآیند آن شامل
‫\لر{Host: sina.sharif.edu} و \لر{URI: /\textasciitilde{}azad/} است. درخواست تصویر
‫\لر{image002.jpg} نیز در فریم \لر{3579} مشاهده شد.
‫
‫\شروع{شکل}[H]
‫\centerimg{4.png}{14cm}
‫\شرح{درخواست‌های \لر{HTTP GET} ثبت‌شده برای صفحه و تصویر آن}
‫\برچسب{شکل:درخواست‌های-get}
‫\پایان{شکل}
‫
‫\section{نتایج بخش \لر{HTTP} و پاسخ پرسش‌ها}
‫
‫\subsection{پروتکل‌های دارای بیشترین حجم ترافیک}
‫
‫آمار سلسله‌مراتب پروتکل‌ها\پانویس{Protocol Hierarchy} بدون اعمال
‫فیلتر نمایشی\پانویس{Display Filter} استخراج شد. راهنمای وایرشارک توضیح می‌دهد
‫که هر ردیف این پنجره تعداد و درصد بسته‌ها و بایت‌های یک پروتکل را نمایش می‌دهد و چون هر بسته
‫چند لایه پروتکلی دارد، مجموع درصد ردیف‌های مختلف الزاما \لر{100} درصد نیست.
‫شکل~\لر{\رجوع{شکل:سلسله‌مراتب}} پنجره آماری آزمایش و جدول~\لر{\رجوع{جدول:آمار-پروتکل‌ها}}
‫مقادیر اصلی استخراج‌شده را نشان می‌دهند.
‫
‫\شروع{شکل}[H]
‫\centerimg{5.png}{14cm}
‫\شرح{آمار سلسله‌مراتب پروتکل‌های موجود در \لر{http.pcap}}
‫\برچسب{شکل:سلسله‌مراتب}
‫\پایان{شکل}
‫
‫\شروع{لوح}[ht]
‫\تنظیم‌ازوسط
‫\شرح{پروتکل‌های اصلی از نظر تعداد بسته و حجم بایت}
‫\شروع{جدول}{|c|c|c|c|c|}
‫\خط‌پر
‫\سیاه پروتکل & \سیاه تعداد بسته & \سیاه درصد بسته‌ها & \سیاه تعداد بایت & \سیاه درصد بایت‌ها \\
‫\خط‌پر \خط‌پر
‫\لر{IPv4} & \لر{5945} & \لر{99.866} & \لر{118900} & \لر{5.731} \\
‫\لر{TCP} & \لر{5716} & \لر{96.019} & \لر{183696} & \لر{8.853} \\
‫\لر{TLS} & \لر{2508} & \لر{42.130} & \لر{1452407} & \لر{70.000} \\
‫\لر{HTTP} & \لر{10} & \لر{0.168} & \لر{3273} & \لر{0.158} \\
‫داده متنی بازسازی‌شده & \لر{2} & \لر{0.034} & \لر{1132648} & \لر{54.589} \\
‫\لر{JPEG} & \لر{2} & \لر{0.034} & \لر{52012} & \لر{2.507} \\
‫\لر{DNS} & \لر{52} & \لر{0.874} & \لر{4112} & \لر{0.198} \\
‫\خط‌پر
‫\پایان{جدول}
‫\برچسب{جدول:آمار-پروتکل‌ها}
‫\پایان{لوح}
‫
‫بنابراین از نظر تعداد بسته، \لر{IPv4} و \لر{TCP} و پس از آن \لر{TLS} غالب‌اند؛ از نظر
‫بایت‌های منتسب به پروتکل‌ها نیز \لر{TLS} با \لر{1452407} بایت و \لر{70} درصد بیشترین سهم را دارد.
‫سهم بالای داده متنی، نتیجه بازسازی بدنه بزرگ سند \لر{HTTP} در وایرشارک است و نباید با جمع
‫ساده درصد لایه‌های مستقل تفسیر شود.
‫
‫\subsection{زمان پاسخ \لر{HTTP} و \لر{Absolute Sequence Number} در \لر{TCP}}
‫
‫نخستین درخواست مرتبط با صفحه اصلی در فریم \لر{3541} و پاسخ \لر{HTTP/1.1 200 OK} متناظر
‫در فریم \لر{3574} قرار دارد. مرجع فیلترهای وایرشارک، فیلد \لر{http.time} را «زمان سپری‌شده از
‫درخواست» تعریف می‌کند. شکل~\لر{\رجوع{شکل:زمان-پاسخ}} مقدار
‫\لر{Time since request = 0.072006025 s} را نشان می‌دهد؛ در نتیجه اختلاف زمانی برابر
‫\لر{72.006025} میلی‌ثانیه است.
‫
‫\شروع{شکل}[H]
‫\centerimg{6.png}{14cm}
‫\شرح{ارتباط فریم درخواست \لر{3541} با پاسخ \لر{3574} و زمان پاسخ اندازه‌گیری‌شده (این تصویر خیلی قابل برش دادن نبود و به این دلیل باید با همین اندازه بزرگ قرار داده می‌شد.)}
‫\برچسب{شکل:زمان-پاسخ}
‫\پایان{شکل}
‫
‫برای بررسی \لر{Absolute Sequence Number}، در \لر{Wireshark} فیلد
‫\لر{Raw Sequence Number} \لر{32}-بیتی آن را نمایش می‌دهد.
‫برای نخستین اتصال \لر{TCP} مرتبط با سرور \لر{HTTP} آزمایش، بسته \لر{SYN} در فریم \لر{3538}
‫از \لر{Port 52332} به \لر{Port 80} ارسال شده و شماره ترتیب خام آن \لر{2987078498} است؛ این مقدار در
‫شکل~\لر{\رجوع{شکل:شماره-ترتیب}} دیده می‌شود.
‫
‫\شروع{شکل}[H]
‫\centerimg{7.png}{7cm}
‫\شرح{شماره ترتیب خام اتصال \لر{TCP} مرتبط با صفحه هدف}
‫\برچسب{شکل:شماره-ترتیب}
‫\پایان{شکل}
‫
‫\subsection{نوع پرس‌وجوها و پاسخ‌های \لر{DNS}}
‫
‫مطابق شکل~\لر{\رجوع{شکل:پرس‌وجوی-dns}}، یکی از پرس‌وجوهای\پانویس{Queries} دامنه
‫\لر{sina.sharif.edu} دارای نوع \لر{A}، کد \لر{1} و کلاس \لر{IN} است. \لر{RFC 1035} نوع
‫\لر{A} را نشانی میزبان و نوع \لر{SOA} با کد \لر{6} را آغاز ناحیه مرجع تعریف می‌کند.
‫
‫\شروع{شکل}[H]
‫\centerimg{8.png}{9cm}
‫\شرح{پرس‌وجوی \لر{DNS} نوع \لر{A} برای \لر{sina.sharif.edu}}
‫\برچسب{شکل:پرس‌وجوی-dns}
‫\پایان{شکل}
‫
‫تحلیل فایل ضبط‌شده نشان داد که مرورگر دو نوع پرس‌وجوی \لر{A} و \لر{HTTPS} ارسال کرده است.
‫پاسخ پرس‌وجوی \لر{A} یک رکورد \لر{A} با مقدار \لر{152.89.13.54} داشت. 
‫همچنین در بخش \لر{Additional} درخواست‌ها و پاسخ‌ها رکورد \لر{OPT} نوع \لر{41} دیده شد. این رکورد
‫طبق \لر{RFC 6891} یک شبه‌رکورد \لر{EDNS(0)} برای اطلاعات کنترلی است و داده پاسخ نام دامنه محسوب
‫نمی‌شود.
‫
‫\subsection{بازیابی تصویر از جریان \لر{HTTP}}
‫
‫\لر{Export Objects}
‫جریان‌های پروتکل\پانویس{Protocol Streams} انتخابی را می‌کاود و
‫\لر{Reconstructed Objects}، از جمله تصویرهای \لر{HTTP}، را ذخیره می‌کند.
‫در شکل~\لر{\رجوع{شکل:فهرست-اشیا}} چهار شیء دیده می‌شود: دو سند \لر{HTML} و دو انتقال از فایل
‫\لر{image002.jpg} در فریم‌های پاسخ \لر{3611} و \لر{4691}. هر ردیف تصویر حدود \لر{26} کیلوبایت است؛ تکرار
‫آن ناشی از بازخوانی صفحه است و تنها یک تصویر متمایز در ضبط وجود دارد.
‫
‫\شروع{شکل}[H]
‫\centerimg{9.png}{12cm}
‫\شرح{فهرست اشیای \لر{HTTP} و پیش‌نمایش \لر{image002.jpg}}
‫\برچسب{شکل:فهرست-اشیا}
‫\پایان{شکل}
‫
‫فایل استخراج‌شده باز و صحت آن بررسی شد. نتیجه در
‫شکل~\لر{\رجوع{شکل:تصویر-بازیابی‌شده}} نمایش داده شده است.
‫
‫\شروع{شکل}[H]
‫\centerimg{10.png}{12cm}
‫\شرح{تصویر \لر{image002.jpg} بازیابی‌شده از جریان \لر{HTTP}}
‫\برچسب{شکل:تصویر-بازیابی‌شده}
‫\پایان{شکل}
‫
‫\section{بررسی ارتباط از طریق \لر{Telnet}}
‫
‫\subsection{مشاهدات نشست زنده \لر{Telehack}}
‫
‫در ضبط زنده، \لر{Client} با نشانی \لر{172.20.10.3} به سرور \لر{Telehack} با نشانی
‫\لر{64.13.139.230} متصل شد. بسته آغاز اتصال، مقدار \لر{Destination Port} را در سرآیند
‫\لر{TCP} برابر \لر{23} نشان می‌دهد. این نتیجه با رجیستری رسمی \لر{IANA} نیز سازگار است که
‫خدمت \لر{Telnet} را برای \لر{TCP Port 23} ثبت کرده است. عبارت
‫\لر{Connected to TELEHACK port 161} در خروجی \لر{Terminal} بخشی از پیام کاربردی سرویس
‫است و نباید با \لر{Destination Port} موجود در سرآیند \لر{TCP} اشتباه شود.
‫
‫جزئیات \لر{TCP Three-way Handshake} نیز در
‫شکل~\لر{\رجوع{شکل:مذاکره-تلنت}} قابل مشاهده است. فریم
‫\لر{21295} یک بسته \لر{SYN} از \لر{172.20.10.3:14580} به \لر{64.13.139.230:23}، فریم
‫\لر{21298} پاسخ \لر{SYN, ACK} سرور و فریم \لر{21299} تأیید نهایی \لر{Client} است. بنابراین
‫نقش دو میزبان فقط از جهت برقراری اتصال و شماره \لر{Port} نیز قابل تشخیص است: آغازگر،
‫\لر{Client} و میزبان پاسخ‌دهنده از \لر{Port 23}، سرور است.
‫
‫در نشست زنده، فرمان‌های \لر{help}، \لر{date} و \لر{cal} اجرا شدند. خروجی شامل منوی متنی
‫\لر{ASCII}، فهرست فرمان‌های موجود، تاریخ و ساعت سامانه و تقویم بود. شکل‌های
‫\لر{\رجوع{شکل:ترمینال-تله‌هک-یک}} و \لر{\رجوع{شکل:ترمینال-تله‌هک-دو}} اجرای این فرمان‌ها و
‫خروجی خوانای آن‌ها را نشان می‌دهند. فرمان \لر{help} نام و توضیح فرمان‌های قابل استفاده را چاپ
‫کرد؛ فرمان \لر{date} مقدار \لر{Monday, August 10, 2026 9:53:15 AM +0330} را برگرداند و
‫فرمان \لر{cal} تقویم ماه‌های \لر{July 2026}، \لر{August 2026} و \لر{September 2026} را
‫نمایش داد.
‫
‫\شروع{شکل}[H]
‫\centerimg{11-telehack-terminal-1.PNG}{9cm}
‫\شرح{اتصال به \لر{Telehack} و اجرای فرمان \لر{help}}
‫\برچسب{شکل:ترمینال-تله‌هک-یک}
‫\پایان{شکل}
‫
‫\شروع{شکل}[H]
‫\centerimg{11-telehack-terminal-2.PNG}{9cm}
‫\شرح{ادامه خروجی \لر{help} و اجرای فرمان‌های \لر{date} و \لر{cal}}
‫\برچسب{شکل:ترمینال-تله‌هک-دو}
‫\پایان{شکل}
‫
‫\subsection{\لر{Telnet Option Negotiation}}
‫
‫بر اساس \لر{RFC 854}، نشست \لر{Telnet} یک اتصال \لر{TCP} است که داده و اطلاعات کنترلی
‫\لر{Telnet} را حمل می‌کند و گزینه‌ها با ساختار \لر{DO}، \لر{DON'T}، \لر{WILL} و
‫\لر{WON'T} مذاکره\پانویس{Negotiate} می‌شوند. \لر{RFC 855} نیز
‫\لر{Subnegotiation} را تبادل پارامترهای یک گزینه
‫پس از پذیرش آن تعریف می‌کند. پیام‌های مشاهده‌شده در ضبط زنده عبارت بودند از:
‫
‫\شروع{فقرات}
‫\فقره \لر{Will Suppress Go Ahead}
‫\فقره \لر{Suboption Terminal Type}
‫\فقره \لر{Do Suppress Go Ahead}
‫\فقره \لر{Do Echo}
‫\فقره \لر{Will Terminal Type}
‫\فقره \لر{Will Negotiate About Window Size}
‫\فقره \لر{Suboption Negotiate About Window Size}
‫\فقره \لر{Won't Environment Option}
‫\فقره \لر{Will New Environment Option}
‫\فقره \لر{Will Binary Transmission}
‫\فقره \لر{Suboption New Environment Option}
‫\پایان{فقرات}
‫
‫هر یک از این پیام‌ها بخشی از تنظیم رفتار \لر{Terminal} است. \لر{Will Suppress Go Ahead} پیشنهاد توقف
‫ارسال پیام کنترلی \لر{GA} و \لر{Do Suppress Go Ahead} درخواست همین رفتار از طرف مقابل است؛
‫این گزینه در \لر{RFC 858} تعریف شده است. پیام \لر{Do Echo} از طرف مقابل می‌خواهد
‫برای نویسه‌های دریافتی \لر{Echo} انجام دهد و معنای آن در \لر{RFC 857} مشخص شده است.
‫
‫در گزینه \لر{Terminal Type}، پیام \لر{Will Terminal Type} آمادگی فرستنده برای اعلام این نوع را نشان می‌دهد
‫و \لر{Suboption Terminal Type} مقدار واقعی آن را منتقل می‌کند.
‫در گزینه اندازه پنجره\پانویس{Window Size} نیز \لر{Will Negotiate About Window Size} پشتیبانی از \لر{NAWS} را
‫اعلام می‌کند و \لر{Subnegotiation} متناظر، عرض و ارتفاع پنجره را از \لر{Client} به سرور می‌فرستد.
‫
‫پیام \لر{Will Binary Transmission} پیشنهاد می‌کند داده‌های آن جهت به‌صورت هشت‌بیتی تفسیر
‫شوند؛ حالت دودویی\پانویس{Binary Mode} برای هر جهت جریان جداگانه مذاکره می‌شود. در نهایت،
‫\لر{Won't Environment Option} رد گزینه قدیمی محیط است، در حالی که
‫\لر{Will New Environment Option} آمادگی برای ارسال متغیرهای محیطی\پانویس{Environment Variables}
‫با سازوکار جدید را اعلام
‫می‌کند و \لر{Suboption New Environment Option} برای درخواست یا انتقال این مقادیر به کار
‫می‌رود.
‫
‫فیلتر \لر{telnet || tcp.port == 23}، اتصال به \لر{Port 23} و این پیام‌های مذاکره را آشکار کرد؛
‫همان‌گونه که در شکل~\لر{\رجوع{شکل:مذاکره-تلنت}} دیده می‌شود، نشانی‌های \لر{Client} و سرور نیز در
‫دو جهت بسته‌ها قابل تشخیص‌اند.
‫
‫\شروع{شکل}[H]
‫\centerimg{12-telehack-negotiation.PNG}{14cm}
‫\شرح{اتصال \لر{TCP} روی \لر{Port 23} و \لر{Telnet Option Negotiation}}
‫\برچسب{شکل:مذاکره-تلنت}
‫\پایان{شکل}
‫
‫پنجره \لر{Follow TCP Stream} نشان داد که هم فرمان‌های \لر{Client} و هم پاسخ‌های سرور قابل خواندن
‫هستند. راهنمای وایرشارک بیان می‌کند که در این پنجره، داده \لر{Client} به سرور قرمز و داده سرور به
‫\لر{Client} آبی نمایش داده می‌شود. در شکل~\لر{\رجوع{شکل:جریان-زنده-تلنت}}، پاسخ‌های آبی
‫سرور شامل پیام خوشامد، منوی \لر{Terminal} و فهرست فرمان‌ها است و داده قرمز، ورودی \لر{Client} را نشان می‌دهد.
‫
‫\شروع{شکل}[H]
‫\centerimg{13-telehack-follow-stream.PNG}{14cm}
‫\شرح{خوانایی داده‌های دو جهت نشست زنده در \لر{Follow TCP Stream}}
‫\برچسب{شکل:جریان-زنده-تلنت}
‫\پایان{شکل}
‫
‫در پایین همین پنجره، جریان \لر{159} با مجموع \لر{7750} بایت نمایش داده شده است. آمار قابل مشاهده
‫شامل \لر{38} بسته \لر{Client}، \لر{66} بسته سرور و \لر{64} نوبت تبادل است. این اعداد فقط به جریان
‫انتخاب‌شده مربوط‌اند. وجود متن آبی خوانا اثبات می‌کند که پاسخ سرور نیز مانند ورودی \لر{Client} از فایل
‫ضبط قابل بازسازی است؛ بنابراین مشاهده به فرمان‌های ارسالی محدود نبود.
‫
‫\subsection{تحلیل شواهد فایل \لر{telnet.pcap}}
‫
‫\subsubsection{پرسش اول: نشانی \لر{IP} مربوط به \لر{Client} و سرور چیست؟}
‫
‫در نخستین بسته \لر{SYN}، مبدأ \لر{192.168.0.2} و مقصد \لر{192.168.0.1} است. همان بسته
‫\لر{Source Port} مربوط به \لر{Client} را \لر{1550} و \لر{Destination Port} سرور را \لر{23} نشان می‌دهد.
‫شکل~\لر{\رجوع{شکل:تلنت-کارخواه-سرور}} این فیلدهای \لر{IPv4} و \لر{TCP} را نمایش می‌دهد.
‫این بسته فریم \لر{1} فایل و تنها نتیجه فیلتر
‫\لر{tcp.flags.syn == 1 \&\& tcp.flags.ack == 0 \&\& tcp.dstport == 23} است. پس
‫\لر{Client} برابر \لر{192.168.0.2} و سرور برابر \لر{192.168.0.1} است.
‫
‫\شروع{شکل}[H]
‫\centerimg{14-telnet-pcap-client-server.PNG}{14cm}
‫\شرح{شناسایی \لر{Client} و سرور \لر{telnet.pcap} از نخستین بسته \لر{SYN}}
‫\برچسب{شکل:تلنت-کارخواه-سرور}
‫\پایان{شکل}
‫
‫\subsubsection{پرسش دوم: نام ورود و رمز عبور استفاده‌شده چیست؟}
‫
‫در نمایش \لر{ASCII} جریان، پس از اعلان \لر{login} مقدار \لر{fake} و پس از اعلان
‫\لر{Password} مقدار \لر{user} در جهت \لر{Client} دیده می‌شود. بنابراین نام ورود \لر{fake} و رمز
‫عبور \لر{user} است. شکل~\لر{\رجوع{شکل:تلنت-رمز}} این داده‌ها و آغاز پاسخ سامانه
‫\لر{OpenBSD/i386} را نشان می‌دهد. پنجره جریان شماره \لر{0} برای کل مکالمه \لر{1634} بایت،
‫\لر{16} بسته \لر{Client}، \لر{30} بسته سرور و \لر{30} نوبت تبادل را گزارش می‌کند. دیده‌شدن مقدار
‫\لر{user} در جهت قرمز نشان می‌دهد که مخفی‌بودن نویسه‌های رمز عبور روی \لر{Terminal} به معنی رمزگذاری
‫آن‌ها در شبکه نیست.
‫
‫\شروع{شکل}[H]
‫\centerimg{15-telnet-pcap-password.PNG}{14cm}
‫\شرح{نام ورود و رمز عبور خوانا در جریان \لر{telnet.pcap}}
‫\برچسب{شکل:تلنت-رمز}
‫\پایان{شکل}
‫
‫\subsubsection{پرسش سوم: چه دستورهایی توسط \لر{Client} اجرا شده‌اند؟}
‫
‫دستورهای \لر{Client} به ترتیب \لر{/sbin/ping www.yahoo.com}، \لر{ls}، \لر{ls -a} و \لر{exit}
‫بودند. فرمان \لر{ping} نام \لر{www.yahoo.com} را در این ضبط به \لر{204.71.200.67} نگاشت و
‫شش پاسخ \لر{ICMP} دریافت کرد. آمار نهایی \لر{6} بسته ارسالی، \لر{6} بسته دریافتی،
‫\لر{0\%} اتلاف و زمان رفت‌وبرگشت کمینه/میانگین/بیشینه
‫\لر{69.885/72.429/75.068 ms} را نشان داد. فرمان \لر{ls} خروجی قابل مشاهده‌ای نداشت، اما
‫\لر{ls -a} پرونده‌های مخفی از جمله \لر{.cshrc}، \لر{.login}، \لر{.mailrc}،
‫\لر{.profile} و \لر{.rhosts} را نمایش داد. سپس \لر{exit} برای پایان نشست ارسال شد. ترتیب کامل
‫در شکل~\لر{\رجوع{شکل:تلنت-دستورها}} دیده می‌شود.
‫
‫\شروع{شکل}[H]
‫\centerimg{16-telnet-pcap-commands.PNG}{14cm}
‫\شرح{ترتیب دستورهای \لر{Client} و پاسخ‌های خوانای سرور در \لر{telnet.pcap}}
‫\برچسب{شکل:تلنت-دستورها}
‫\پایان{شکل}
‫
‫مشاهده مستقیم نام ورود، رمز عبور، دستورها و پاسخ‌های سرور در حالت \لر{ASCII} نشان می‌دهد که
‫محتوای این نشست بدون محرمانگی\پانویس{Confidentiality} در فایل ضبط قابل بازیابی است. در نشست زنده، فرمان‌های
‫\لر{help}، \لر{date} و \لر{cal} و خروجی آن‌ها خوانا بودند؛ در فایل تأمین‌شده نیز نام ورود
‫\لر{fake}، رمز عبور \لر{user} و چهار فرمان \لر{Client} آشکار شدند. بنابراین تحلیل فقط به
‫سرآیندهای \لر{TCP} محدود نیست و محتوای کاربردی نشست نیز قابل مشاهده است. این نتیجه تجربی خطر
‫استفاده از \لر{Telnet} برای انتقال اطلاعات حساس را نشان می‌دهد.
‫
‫\section{بررسی درخواست و پاسخ‌های \لر{DNS}}
‫
‫\subsection{آماده‌سازی و تنظیمات شبکه}
‫
‫آزمایش سوم روی \لر{Wireless LAN adapter Wi-Fi} در ویندوز انجام شد. پیش از ایجاد پرس‌وجو،
‫دستور \لر{ipconfig /flushdns} اجرا شد و پیام موفقیت پاک‌شدن حافظه نهان\پانویس{Cache}
‫مفسر\پانویس{Resolver}
‫\لر{DNS} نمایش داده شد. شکل~\لر{\رجوع{شکل:پاک‌سازی-dns}} اجرای موفق دستور را
‫نشان می‌دهد.
‫
‫\شروع{شکل}[H]
‫\centerimg{17-dns-resolver-settings-1.PNG}{10cm}
‫\شرح{پاک‌سازی موفق حافظه نهان مفسر \لر{DNS} در ویندوز}
‫\برچسب{شکل:پاک‌سازی-dns}
‫\پایان{شکل}
‫
‫خروجی \لر{ipconfig /all} در شکل~\لر{\رجوع{شکل:تنظیمات-dns}}، \لر{Interface} فعال و تنظیمات آن را
‫نشان می‌دهد. نشانی \لر{IPv4} برابر \لر{172.20.10.3} و \لر{IPv6 Link-local Address}
‫برابر \لر{fe80::da47:c186:397:ada3\%9} است. \لر{Upstream DNS Servers}
‫پیکربندی‌شده عبارت‌اند از
‫\لر{fe80::ac00:7aff:fea4:3d64\%9} و \لر{172.20.10.1}.
‫
‫\شروع{شکل}[H]
‫\centerimg{17-dns-resolver-settings-2.PNG}{12cm}
‫\شرح{نشانی‌های \لر{Wi-Fi Interface} و سرورهای \لر{DNS} پیکربندی‌شده}
‫\برچسب{شکل:تنظیمات-dns}
‫\پایان{شکل}
‫
‫پسوند \لر{\%9} در خروجی ویندوز شناسه \لر{Interface} محلی است و همان نشانی در ستون‌های وایرشارک بدون
‫این پسوند دیده می‌شود. اجرای \لر{/flushdns} تنها وجود
‫حافظه نهان سرویس \لر{DNS Client} ویندوز را نشان می‌دهد و آن را به مقصد شبکه تبدیل نمی‌کند.
‫
‫فرمان \لر{nslookup -type=A sina.sharif.edu} اجرا شد. همان‌گونه که در
‫شکل~\لر{\رجوع{شکل:خروجی-nslookup}} دیده می‌شود، خروجی ابتدا \لر{Server: Unknown} و
‫\لر{Address: 8.8.8.8} را نمایش داد و دو \لر{Timeout} دوثانیه‌ای را ثبت کرد، اما در
‫نهایت یک \لر{Non-authoritative Answer} شامل \لر{152.89.13.54} برگرداند.
‫وقتی آرگومان دوم \لر{nslookup} داده نشود، برنامه از سرور پیش‌فرض
‫استفاده می‌کند.
‫
‫\شروع{شکل}[H]
‫\centerimg{18-dns-dig-output.PNG}{10cm}
‫\شرح{اجرای \لر{nslookup} و دریافت پاسخ \لر{A} برای \لر{sina.sharif.edu}}
‫\برچسب{شکل:خروجی-nslookup}
‫\پایان{شکل}
‫
‫\subsection{پرسش اول: درخواست به کدام سرور ارسال و پاسخ از کدام سرور دریافت شد؟}
‫
‫برای جداسازی زوج مورد نظر، فیلترهای
‫\لر{dns.flags.response == 0 \&\& dns.qry.name == "sina.sharif.edu"} و
‫\لر{dns.flags.response == 1 \&\& dns.qry.name == "sina.sharif.edu"} استفاده شدند. فیلدهای \لر{dns.flags.response} و \لر{dns.qry.name} برای تشخیص جهت پیام و
‫نام پرسش تعریف می‌شوند.
‫
‫در فریم \لر{1326}، درخواست از نشانی محلی \لر{fe80::da47:c186:397:ada3} به
‫\لر{fe80::ac00:7aff:fea4:3d64} ارسال شد. پاسخ متناظر در فریم \لر{1337} از
‫\لر{fe80::ac00:7aff:fea4:3d64} به \لر{fe80::da47:c186:397:ada3} بازگردانده شد. بنابراین مقصد درخواست و مبدأ پاسخ هر دو \لر{fe80::ac00:7aff:fea4:3d64} هستند. این
‫نشانی با \لر{Link-local DNS Server} موجود در \لر{ipconfig /all} یکسان است؛ فقط شناسه \لر{Interface}
‫\لر{\%9} در نمایش بسته وجود ندارد.
‫
‫شکل~\لر{\رجوع{شکل:زوج-dns}} هر دو فریم را در فهرست بسته‌ها نشان می‌دهد. درخواست با
‫\لر{UDP} از \لر{Port 59746} به \لر{Port 53} و پاسخ در جهت معکوس از \لر{Port 53} به
‫\لر{Port 59746} منتقل شد. \لر{Transaction ID} مشترک \لر{0x0002} نیز ارتباط دو
‫پیام را تأیید می‌کند.
‫
‫\شروع{شکل}[H]
‫\centerimg{19-dns-request-response-list.PNG}{14cm}
‫\شرح{درخواست فریم \لر{1326} و پاسخ فریم \لر{1337} با شناسه \لر{0x0002}}
‫\برچسب{شکل:زوج-dns}
‫\پایان{شکل}
‫
‫\subsection{پرسش دوم: تحلیل سرآیندهای درخواست و پاسخ \لر{DNS}}
‫
‫بر اساس \لر{RFC 1035}، سرآیند \لر{DNS} شامل شناسه، پرچم‌ها\پانویس{Flags} و چهار شمارنده
‫بخش‌های \لر{Question}، \لر{Answer}، \لر{Authority} و \لر{Additional} است.
‫در ادامه، هر دو پیام مطابق همین ساختار تحلیل می‌شوند.
‫
‫\subsubsection{سرآیند درخواست}
‫
‫شکل‌های~\لر{\رجوع{شکل:سرآیند-درخواست-dns-یک}} و
‫\لر{\رجوع{شکل:سرآیند-درخواست-dns-دو}} لایه‌های \لر{IPv6}، \لر{UDP} و \لر{DNS} درخواست را
‫نشان می‌دهند. فریم \لر{1326} دارای اندازه \لر{95} بایت است و \لر{UDP} را با
‫\لر{Source Port 59746}، \لر{Destination Port 53}، طول \لر{41} بایت و بار مفید \لر{DNS} به طول \لر{33} بایت
‫حمل می‌کند.
‫
‫\شروع{شکل}[H]
‫\centerimg{20-dns-request-header-1.PNG}{12cm}
‫\شرح{نشانی‌های \لر{IPv6} و \لر{UDP Ports} درخواست \لر{DNS}}
‫\برچسب{شکل:سرآیند-درخواست-dns-یک}
‫\پایان{شکل}
‫
‫\شروع{شکل}[H]
‫\centerimg{20-dns-request-header-2.PNG}{12cm}
‫\شرح{پرچم‌ها، شمارنده‌ها و بخش پرسش فریم \لر{1326}}
‫\برچسب{شکل:سرآیند-درخواست-dns-دو}
‫\پایان{شکل}
‫
‫\لر{Transaction ID} برابر \لر{0x0002} و مقدار پرچم‌ها \لر{0x0100} است. بیت \لر{QR} صفر است، پس پیام
‫یک درخواست است؛ \لر{Opcode} صفر به معنی پرس‌وجوی استاندارد، \لر{TC=0} به معنی قطع‌نشدن
‫پیام و \لر{RD=1} به معنی درخواست بازگشتی\پانویس{Recursive Query} است. همچنین \لر{Z=0}
‫بیت رزروشده، \لر{AD=0} نبود داده احرازشده و \لر{CD=0} نپذیرفتن داده احرازنشده را نشان
‫می‌دهند. شمارنده‌ها شامل
‫\لر{Questions=1}، \لر{Answer RRs=0}، \لر{Authority RRs=0} و \لر{Additional RRs=0}
‫هستند. تنها پرسش برای نام \لر{sina.sharif.edu}، نوع \لر{A} با کد \لر{1} و کلاس \لر{IN}
‫با کد \لر{0x0001} است. وایرشارک همچنین \لر{Response In: 1337} را برای این درخواست نشان
‫می‌دهد.
‫
‫\subsubsection{سرآیند پاسخ}
‫
‫شکل‌های~\لر{\رجوع{شکل:سرآیند-پاسخ-dns-یک}} و
‫\لر{\رجوع{شکل:سرآیند-پاسخ-dns-دو}} پاسخ فریم \لر{1337} را نشان می‌دهند. این فریم
‫\لر{111} بایت است و با \لر{UDP} از \لر{Port 53} سرور به \لر{Port 59746} در \لر{Client} بازمی‌گردد.
‫
‫\شروع{شکل}[H]
‫\centerimg{21-dns-response-header-1.PNG}{14cm}
‫\شرح{نشانی‌ها، \لر{Ports} و آغاز سرآیند پاسخ \لر{DNS}}
‫\برچسب{شکل:سرآیند-پاسخ-dns-یک}
‫\پایان{شکل}
‫
‫\شروع{شکل}[H]
‫\centerimg{21-dns-response-header-2.PNG}{14cm}
‫\شرح{پرچم‌ها و شمارنده‌های پاسخ فریم \لر{1337}}
‫\برچسب{شکل:سرآیند-پاسخ-dns-دو}
‫\پایان{شکل}
‫
‫\لر{Transaction ID} پاسخ نیز \لر{0x0002} و پرچم‌های آن \لر{0x8180} است. \لر{QR=1} پاسخ‌بودن
‫پیام، \لر{Opcode=0} پرس‌وجوی استاندارد، \لر{AA=0} نشان‌دهنده \لر{Non-authoritative} بودن پاسخ،
‫\لر{TC=0} کامل‌بودن پیام، \لر{RD=1} باقی‌ماندن درخواست
‫بازگشتی، \لر{RA=1} امکان انجام بازگشت توسط سرور\پانویس{Recursion Available} و
‫\لر{Z=0} بیت رزروشده، \لر{AD=0} نبود داده احرازشده و \لر{CD=0} نپذیرفتن داده احرازنشده
‫را مشخص می‌کنند؛ \لر{RCODE=0} نیز به معنی نبود خطا است. شمارنده‌ها برابر \لر{Questions=1}،
‫\لر{Answer RRs=1}، \لر{Authority RRs=0} و \لر{Additional RRs=0} هستند. تفاوت شمارنده
‫پاسخ با درخواست، افزوده‌شدن یک رکورد منبع\پانویس{Resource Record} در بخش \لر{Answer} است.
‫
‫\subsubsection{رکورد پاسخ و نتیجه تحلیل}
‫
‫شکل~\لر{\رجوع{شکل:رکورد-پاسخ-dns}} پرسش و رکورد پاسخ را نمایش می‌دهد. نام، نوع و کلاس پرسش
‫\لر{sina.sharif.edu}، \لر{A} و \لر{IN} است. پاسخ نیز یک رکورد \لر{A} از کلاس \لر{IN} با
‫زمان ماندگاری\پانویس{Time To Live} \لر{77 s}، طول داده \لر{4} بایت و مقدار
‫\لر{152.89.13.54} دارد. مقدار \لر{Request In: 1326} و زمان پاسخ
‫\لر{163.464800 ms} نیز در همین تصویر ثبت شده است.
‫
‫\شروع{شکل}[H]
‫\centerimg{22-dns-answer-records-1.PNG}{14cm}
‫\شرح{رکورد \لر{A} پاسخ، زمان ماندگاری و زمان پاسخ اندازه‌گیری‌شده}
‫\برچسب{شکل:رکورد-پاسخ-dns}
‫\پایان{شکل}
‫
‫برابری شناسه \لر{0x0002} دو پیام و ارجاع‌های \لر{Response In} و \لر{Request In} آن‌ها را
‫به یک تراکنش متصل می‌کند. تغییر \لر{QR} از صفر به یک، پاسخ‌شدن پیام را نشان می‌دهد؛ باقی‌ماندن
‫\لر{RD} و فعال‌شدن \لر{RA} یعنی \لر{Client} بازگشت را درخواست کرده و سرور امکان آن را اعلام کرده
‫است. \لر{RCODE=0} و \لر{Answer RRs=1} دریافت موفق پاسخ را تأیید می‌کنند و \لر{AA=0} با
‫عبارت \لر{Non-authoritative answer} در خروجی \لر{nslookup} سازگار است.
‫
‫\section{جمع‌بندی}
‫
‫در این آزمایش، درخواست‌های خوانای \لر{HTTP} برای میزبان \لر{sina.sharif.edu} ثبت و
‫پاسخ متناظر نخستین درخواست در \لر{72.006025} میلی‌ثانیه دریافت شد. \لر{TCP} و \لر{IPv4} از
‫نظر تعداد بسته و \لر{TLS} از نظر حجم بایت بر ضبط غالب بودند. شماره ترتیب خام اتصال
‫\لر{HTTP} برابر \لر{2987078498} به دست آمد. مرورگر پرس‌وجوهای \لر{DNS} نوع \لر{A} و
‫\لر{HTTPS} صادر کرد و فایل \لر{image002.jpg} نیز از جریان \لر{HTTP} بازیابی شد.
‫در بخش \لر{Telnet}، اتصال زنده \لر{Telehack} میان \لر{172.20.10.3} و
‫\لر{64.13.139.230} روی \لر{TCP Port 23} برقرار شد و فرمان‌های \لر{help}، \لر{date} و
‫\لر{cal} همراه با پاسخ‌های سرور خوانا بودند. شواهد \لر{telnet.pcap} نیز \لر{Client}
‫\لر{192.168.0.2}، سرور \لر{192.168.0.1}، نام ورود \لر{fake}، رمز عبور \لر{user} و ترتیب
‫چهار فرمان \لر{Client} را آشکار کرد.
‫در بخش \لر{DNS}، درخواست \لر{A} دامنه \لر{sina.sharif.edu} در فریم \لر{1326} به سرور
‫\لر{fe80::ac00:7aff:fea4:3d64} ارسال و پاسخ موفق آن در فریم \لر{1337} دریافت شد. پاسخ یک
‫رکورد \لر{A} با زمان ماندگاری \لر{77} ثانیه و مقدار \لر{152.89.13.54} داشت و تحلیل پرچم‌ها
‫نشان داد که پاسخ بدون خطا و بازگشتی است و از نظر سرور مرجع \لر{Non-authoritative} محسوب می‌شود.
‫
‫\section{منابع}
‫
‫\شروع{شمارش}
‫
‫
‫\فقره
‫\href{https://www.rfc-editor.org/rfc/rfc1035.html}{\لر{IETF RFC 1035, ``Domain Names --- Implementation and Specification''}}، تعریف نوع‌های \لر{A} و \لر{SOA} و کلاس \لر{IN}.
‫
‫\فقره
‫\href{https://www.rfc-editor.org/rfc/rfc9460.html}{\لر{IETF RFC 9460, ``Service Binding and Parameter Specification via the DNS''}}، تعریف رکورد \لر{HTTPS} با نوع \لر{65}.
‫
‫\فقره
‫\href{https://www.rfc-editor.org/rfc/rfc6891.html}{\لر{IETF RFC 6891, ``Extension Mechanisms for DNS (EDNS(0))''}}، تعریف شبه‌رکورد \لر{OPT} با نوع \لر{41}.
‫
‫\فقره
‫\href{https://www.rfc-editor.org/rfc/rfc854.html}{\لر{IETF RFC 854, ``Telnet Protocol Specification''}}، تعریف اتصال و \لر{Option Negotiation} در \لر{Telnet}.
‫
‫\فقره
‫\href{https://www.rfc-editor.org/rfc/rfc855.html}{\لر{IETF RFC 855, ``Telnet Option Specifications''}}، تعریف روش \لر{Subnegotiation} گزینه‌های \لر{Telnet}.
‫
‫\فقره
‫\href{https://www.rfc-editor.org/rfc/rfc858.html}{\لر{IETF RFC 858, ``Telnet Suppress Go Ahead Option''}}، تعریف گزینه حذف پیام \لر{Go Ahead}.
‫
‫\فقره
‫\href{https://www.rfc-editor.org/rfc/rfc857.html}{\لر{IETF RFC 857, ``Telnet Echo Option''}}، تعریف گزینه \لر{Echo} برای نویسه‌های دریافتی.
‫
‫\فقره
‫\href{https://www.rfc-editor.org/rfc/rfc1091.html}{\لر{IETF RFC 1091, ``Telnet Terminal-Type Option''}}، تعریف \لر{Negotiation} و \لر{Subnegotiation} گزینه \لر{Terminal Type}.
‫
‫\فقره
‫\href{https://www.rfc-editor.org/rfc/rfc1073.html}{\لر{IETF RFC 1073, ``Telnet Window Size Option''}}، تعریف گزینه \لر{NAWS} و انتقال اندازه پنجره.
‫
‫\فقره
‫\href{https://www.rfc-editor.org/rfc/rfc856.html}{\لر{IETF RFC 856, ``Telnet Binary Transmission''}}، تعریف انتقال دودویی هشت‌بیتی.
‫
‫\فقره
‫\href{https://www.rfc-editor.org/rfc/rfc1572.html}{\لر{IETF RFC 1572, ``Telnet Environment Option''}}، تعریف گزینه جدید انتقال متغیرهای محیطی.
‫
‫\پایان{شمارش}
‫